\section{Capitolo 0: Dal Testo alla Logica}

In questo capitolo preliminare, impariamo a decifrare il "linguaggio da esame". Spesso la difficoltà non è scrivere il codice, ma capire \textit{esattamente} cosa viene chiesto.
Usiamo il \textbf{Dizionario Tecnico} per tradurre le parole chiave in codice.

\subsection{1. Il Glossario delle Strutture Dati}

\begin{dizionario}{Array Annidato (Nested)}
\textbf{Definizione:} Un array che contiene al suo interno altri array.
\textbf{Esempio:} \texttt{[1, [2, 3], 4]}.
\textbf{Cosa Significa nel Codice:} Probabilmente dovrai usare la \textbf{ricorsione} (vedi `flat` o `partition`) o doppi cicli.
\end{dizionario}

\begin{dizionario}{Insieme (Set) implementato come Oggetto}
\textbf{Definizione:} Un oggetto dove la chiave è l'elemento e il valore è \texttt{true}.
\textbf{Codice:} \texttt{\{ "mela": true, "pera": true \}}.
\textbf{Perché?} Per avere un accesso veloce $O(1)$ invece di cercare nell'array $O(n)$.
\end{dizionario}

\begin{dizionario}{Multi-insieme}
\textbf{Definizione:} Come un Set, ma contiamo quante volte appare un elemento.
\textbf{Codice:} \texttt{\{ "mela": 2, "pera": 1 \}}.
\textbf{Pattern Tipico:} \texttt{if (k in obj) obj[k]++; else obj[k] = 1;}
\end{dizionario}

\subsection{2. Il Glossario per le Funzioni}

\begin{dizionario}{Predicato}
\textbf{Definizione:} Una funzione che restituisce \textbf{sempre e solo} \texttt{boolean}.
\textbf{Firma:} \texttt{f: (x: T) => boolean}
\textbf{Esempio:} \texttt{isPari(x)}, \texttt{isMaggiorenne(user)}.
\end{dizionario}

\begin{dizionario}{Funzione Fabbrica (Factory)}
\textbf{Definizione:} Una funzione che \textbf{crea e restituisce} una nuova funzione.
\textbf{Concetto Chiave:} \textbf{Closure}. La funzione restituita "ricorda" i parametri della fabbrica.
\end{dizionario}

\begin{dizionario}{In-Place (Side Effect)}
\textbf{Definizione:} Modificare direttamente l'oggetto/array passato come argomento, SENZA crearne uno nuovo.
\textbf{Attenzione:} Se il testo dice "senza side-effect" o "funzionale pura", devi clonare l'input prima di toccarlo!
\end{dizionario}

\subsection{Esercizi di "Traduzione"}

\begin{esercizio}{Traduzione 1: La "Fabbrica di Predicati"}
\textbf{TESTO:} "Scrivere una funzione `createFilter` che accetta un numero `soglia` e restituisce un predicato che vale true se l'input è maggiore della soglia."

\textbf{ANALISI:}
1. \textit{Scrivere una funzione `createFilter` che accetta `soglia`}: \texttt{function createFilter(soglia) \{ ... \}}
2. \textit{Restituisce un predicato}: Deve fare \texttt{return function(x) \{ ... \}}
3. \textit{Vale true se input > soglia}: \texttt{return x > soglia}

\end{esercizio}

\begin{jsbox}
function createFilter(soglia) {
    // Ritorna una funzione (Closure) che ricorda 'soglia'
    return function(valoreInput) {
        return valoreInput > soglia; // Ritorna boolean (Predicato)
    };
}

const filtroMaggiorenni = createFilter(18);
console.log(filtroMaggiorenni(20)); // true
\end{jsbox}

\begin{esercizio}{Traduzione 2: "Senza Funzioni di Libreria"}
\textbf{TESTO:} "Calcolare l'intersezione tra due array A e B senza usare `Set`, `filter` o `includes`."

\textbf{ANALISI:}
1. \textit{Intersezione}: Elementi comuni.
2. \textit{Senza `includes`}: Devi farti un doppio ciclo for a mano o usare un oggetto di appoggio (Dictionary) per la velocità.
\end{esercizio}

\begin{jsbox}
function intersezioneManuale(A, B) {
    let result = [];
    // Strategia Dizionario per efficienza
    let mappaA = {};
    
    // 1. Popolo il dizionario con gli elementi di A
    for (let i = 0; i < A.length; i++) {
        mappaA[A[i]] = true;
    }
    
    // 2. Controllo se gli elementi di B sono nel dizionario
    for (let j = 0; j < B.length; j++) {
        let val = B[j];
        if (mappaA[val]) { // Equivale a A.includes(val)
            result.push(val);
        }
    }
    return result;
}
\end{jsbox}

\subsection{Tabella Rapida Concetti Matematici}
\begin{enumerate}
    \item \textbf{Suriettiva:} Copre tutto il codominio.
    \item \textbf{Iniettiva:} Nessuna collisione (input diversi $\to$ output diversi).
    \item \textbf{Invertibile:} Biunivoca (1-a-1 perfetto).
    \item \textbf{Punto Fisso:} Quando $f(x) = x$.
\end{enumerate}
